% Source pp. 111-134 (Lectures13-18.pdf).
\section{Systematics of Renormalization}

% local diagram pieces for this section
\newcommand{\legs}[1]{\foreach \p/\a/\n in {#1} {\draw[prop,densely dashed] (\p) -- ++(\a:1.1) node[pos=1.18,font=\tiny] {\n};}}
\newcommand{\bubS}{\diac[0.4]{\coordinate (L) at (-0.9,0); \coordinate (R) at (0.9,0);
  \draw[prop] (0,0) ellipse (0.9 and 0.55); \legs{L/135/1,L/225/2,R/45/3,R/315/4}
  \node[vtx] at (L) {}; \node[vtx] at (R) {};}}
\newcommand{\bubT}{\diac[0.4]{\coordinate (T) at (0,0.8); \coordinate (B) at (0,-0.8);
  \draw[prop] (0,0) ellipse (0.55 and 0.8); \legs{T/150/1,T/30/3,B/210/2,B/330/4}
  \node[vtx] at (T) {}; \node[vtx] at (B) {};}}
\newcommand{\bubU}{\diac[0.4]{\coordinate (T) at (0,0.8); \coordinate (B) at (0,-0.8);
  \draw[prop] (0,0) ellipse (0.55 and 0.8);
  \draw[prop,densely dashed] (T) -- ++(150:1.1) node[pos=1.18,font=\tiny] {1};
  \draw[prop,densely dashed] (B) -- ++(210:1.1) node[pos=1.18,font=\tiny] {2};
  \draw[prop,densely dashed] (T) .. controls (0.9,0.6) and (0.7,-1.4) .. (1.2,-1.6) node[right,font=\tiny] {4};
  \draw[prop,densely dashed] (B) .. controls (0.9,-0.6) and (0.7,1.4) .. (1.2,1.6) node[right,font=\tiny] {3};
  \node[vtx] at (T) {}; \node[vtx] at (B) {};}}
\newcommand{\ctX}[1]{\diac[0.4]{\coordinate (O) at (0,0); \legs{O/135/1,O/225/2,O/45/3,O/315/4}
  \node[otimesv,minimum size=5pt] at (O) {}; \node[font=\tiny] at (0,0.55) {(#1)};}}
\newcommand{\ctline}[1]{\dia[0.4]{\draw[prop,densely dashed] (-2,0) -- (2,0); \node[otimesv,minimum size=5pt] at (0,0) {}; \node[font=\tiny] at (0,0.55) {(#1)};}}

If the theory under consideration is renormalizable, like $\varphi^4$ theory in
4 dimensions, one can develop renormalized perturbation theory. According to the
program outlined above one has to start with the action in its ``bare'' form,
implicitly equipped with some ``regularization'', or ``cutoff'', i.e.\ with some
modification of the theory at the momenta $\ge\Lambda$ which makes all integrals
in perturbation theory convergent. There are many possible implementations of
the cutoff which can be more or less convenient depending on circumstances. Let
us discuss some more frequently used ones.

\subsection*{Regularization methods}

\paragraph{Lattice regularization:} We already discussed this regularization in
relation with the definition of the functional integral. In this approach
continuous space $\R^d$ is replaced by $d$-dimensional, say hypercubic, lattice
with some lattice spacing $\Delta$ which plays the role of inverse cutoff
parameter $\Lambda^{-1}$,
\[
x\to x_n=\Delta\sum_{\mu=1,\dots,d}e_\mu\,n_\mu\,,
\]
where $e_\mu$ is the unit vector in the direction $\mu$ and $n_1,\dots,n_d$ are
integers. The lattice action is obtained by replacing derivatives by finite
differences,
\begin{equation}\label{eq:Alat}
\mathcal A_{\text{lat}}=\Delta^d\sum_{x\in\Delta\mathbb Z^d}\left[\frac12\sum_{\mu=1,\dots,d}\left(\frac{\varphi_0(x+\Delta e_\mu)-\varphi_0(x)}\Delta\right)^2
+\frac{m_0^2}2\,\varphi_0^2(x)+V\bigl(\varphi_0(x)\bigr)\right]
\end{equation}
and the integration measure in the functional integral is taken as
\[
[D\varphi]\to\prod_{x\in\Delta\mathbb Z^d}d\varphi_0(x)\,.
\]
Advantage of this method is that it makes sense non-perturbatively. The lattice
theory \eqref{eq:Alat} is perfectly meaningful in classical statistical
mechanics. Although it is hardly possible to calculate the integral
analytically, one can try to apply numerical approach, say using the
Monte-Carlo method. This approach is known as the lattice Field Theory. At the
same time perturbative calculations on the lattice are relatively complicated.
In addition, continuous space symmetries are broken, and one only expects them
to be restored after the cutoff is removed.

\paragraph{Proper-time regularization:} We have seen before that the
momentum-space propagator $\tilde D(p)$ admits Schwinger's proper-time
representation
\[
\tilde D(p)=\int_0^\infty ds\,e^{-s(p^2+m^2)}\,,
\]
where $s$ is interpreted as the rescaled length of the path of relativistic
particle in the Euclidean space-time. One can exclude the paths which are
``too short'' by replacing $\tilde D(p)\to\tilde D_\Lambda(p)$,
\[
\tilde D_\Lambda(p)=\int_{1/\Lambda^2}^\infty ds\,e^{-s(p^2+m^2)}\,.
\]
Note that this is particular case of the regularized propagator we considered
before,
\begin{equation}\label{eq:regprop}
\frac1{p^2+m^2}\,\Phi(p^2/\Lambda^2)\,,
\end{equation}
with $\Phi(x)=e^{-x}$ (up to the constant factor $e^{-m^2/\Lambda^2}$, which
tends to $1$ as $\Lambda\to\infty$).

\paragraph{Pauli-Villars regularization} is another version of
\eqref{eq:regprop}, with
\begin{equation}\label{eq:PV}
\Phi_{PV}(p^2/\Lambda^2)=\frac{\Lambda^2}{p^2+\Lambda^2}\,.
\end{equation}
In this case the cutoff propagator takes the form
\[
\tilde D(p)\to\frac1{p^2+m^2}\,\Phi_{PV}(p^2/\Lambda^2)\to\frac1{p^2+m^2}-\frac1{p^2+\Lambda^2}\,,
\]
equivalent to \eqref{eq:PV} when $\Lambda\gg m$.

\paragraph{Dimensional regularization:} This is the most technically advanced
(although perhaps the least physically transparent) regularization method. Its
advantage is that it usually preserves important symmetries (in particular the
gauge symmetry in the gauge theories) and significantly simplifies
calculations, and for these reasons this method is used today in perturbative
calculations more frequently than others. The idea is this. We have already
observed that lowering the space dimension generally improves the
large-momentum convergence of Feynman diagrams. Suppose that for $d$
sufficiently low the momentum integral associated with given diagram is
convergent without any additional cutoff. Suppose that in addition we have
managed to calculate this integral as an analytic function of $d$. Then we can
analytically continue the result to the ``physical value'' $d=4$; the
4-dimensional divergences manifest themselves as singularities (poles) in the
variable $d$ at $d=4$. On more formal level one can define the dimensional
continuation of momentum integrals by three conditions, obviously valid for
convergent integrals, and postulated for any $d$:
\[
\begin{array}{lll}
\text{(i)} & \displaystyle\int d^dk\,F(k+p)=\int d^dk\,F(k) & \text{translation}\\[10pt]
\text{(ii)} & \displaystyle\int d^dk\,F(Ck)=|C|^{-d}\int d^dk\,F(k) & \text{dilation}\\[10pt]
\text{(iii)} & \displaystyle\int d^{d_1}k_1\,d^{d_2}k_2\,F(k_1)G(k_2)=\int d^{d_1}k\,F(k_1)\int d^{d_2}k_2\,G(k_2)\,,\quad d=d_1+d_2 & \text{factorization}
\end{array}
\]
Let me illustrate this approach by simple example. Consider the integral
\begin{equation}\label{eq:dimreg1}
\int\frac{d^dk}{(2\pi)^d}\,\frac1{(k^2+m^2)^2}\,.
\end{equation}
This integral converges for $d<4$. For any $d<4$ it can be transformed as
follows
\begin{equation}\label{eq:dimreg2}
\int_0^\infty s\,ds\int\frac{d^dk}{(2\pi)^d}\,e^{-s(k^2+m^2)}=\int_0^\infty s\,ds\,(4\pi s)^{-\frac d2}\,e^{-m^2s}\,,
\end{equation}
where the properties (iii) and (ii) were used in evaluating the $k$-integral.
Using the definition of the Euler's Gamma function
\[
\Gamma(z)=\int_0^\infty dt\,t^{z-1}\,e^{-t}
\]
we obtain for $d<4$
\begin{equation}\label{eq:dimreg3}
\frac{(m^2)^{\frac d2-2}}{(4\pi)^{\frac d2}}\,\Gamma\left(2-\frac d2\right)\,.
\end{equation}
As is known the Gamma function $\Gamma(z)$ is analytic (actually, meromorphic)
function of $z$ with poles at $z=0,-1,-2,\dots$. In particular, the above
expression has pole at $d=4$, i.e.\ exactly where the integral diverges
logarithmically. For $d=4-\epsilon$, $\epsilon\to0$
\[
\Gamma\left(2-\frac d2\right)=\frac2\epsilon-\gamma+O(\epsilon)
\]
where $\gamma=0.5772\ldots$ is the Euler's constant.

It is interesting to explain on this example the interplay between dimensional
continuation and cutoff regularization. If we regularize the propagators in
\eqref{eq:dimreg1}, say by proper-time regularization (which amounts to
replacing $\int_0^\infty ds\to\int_{1/\Lambda^2}^\infty ds$ in
\eqref{eq:dimreg2}) the integral \eqref{eq:dimreg1} becomes regular function of
$d$ for all $d$ including $d=4$. In the neighborhood of $d=4$ one would have
instead of \eqref{eq:dimreg2}
\[
\frac{(m^2)^{\frac d2-2}}{(4\pi)^{\frac d2}}\int_{m^2/\Lambda^2}^\infty dt\,t^{1-\frac d2}\,e^{-t}\,,\qquad t=m^2s\,.
\]
The integral here can be transformed using integration by parts,
\[
\int_{m^2/\Lambda^2}^\infty dt\,t^{1-\frac d2}\,e^{-t}=\frac2{4-d}\Bigl(t^{2-\frac d2}\,e^{-t}\Big|_{m^2/\Lambda^2}^\infty+\int_{m^2/\Lambda^2}^\infty dt\,t^{2-\frac d2}\,e^{-t}\Bigr)\,.
\]
For $\epsilon=4-d\ll1$ and $m^2\ll\Lambda^2$ the cutoff integral becomes
\[
\frac1{(4\pi)^{\frac d2}}\,\frac2{4-d}\,\Bigl((m^2)^{\frac d2-2}-(\Lambda^2)^{\frac d2-2}\Bigr)\,.
\]
If at fixed $d<4$ we send the cutoff momentum $\Lambda$ to infinity we obtain
the ``dimensionally continued'' expression with the pole at $d=4$. On the other
hand, if at fixed $\Lambda$ we take the limit $d\to4$ we get a finite result
with $\log\Lambda^2$ replacing the pole in $4-d$.

With this understanding one can often use dimensional continuation to do the
calculations efficiently and then interpret the results in terms of cutoff
regularization. Instead of giving the bare parameters $m_0^2$, $\lambda_0$ and
$Z$ dependence on the cutoff momentum $\Lambda$, within the framework of
dimensional regularization one adjusts their dependence on $\epsilon=4-d$, so
that the correlation functions of renormalized fields $\varphi$ have finite
limit $d\to4$.

Some generalizations of the integral \eqref{eq:dimreg1} will be useful in
future calculations. For instance,
\begin{equation}\label{eq:dimregalpha}
\int\frac{d^dk}{(2\pi)^d}\,\frac1{(k^2+m^2)^\alpha}=\frac1{\Gamma(\alpha)}\int_0^\infty ds\,s^{\alpha-1}\int\frac{d^dk}{(2\pi)^d}\,e^{-s(k^2+m^2)}
=\frac{(m^2)^{\frac d2-\alpha}}{(4\pi)^{\frac d2}}\,\frac{\Gamma(\alpha-\frac d2)}{\Gamma(\alpha)}\,,
\end{equation}
\begin{equation}\label{eq:dimregk2}
\begin{split}
\int\frac{d^dk}{(2\pi)^d}\,\frac{k^2}{(k^2+m^2)^\alpha}&=\int\frac{d^dk}{(2\pi)^d}\left[\frac1{(k^2+m^2)^{\alpha-1}}-\frac{m^2}{(k^2+m^2)^\alpha}\right]\\
&=\frac{(m^2)^{\frac d2-\alpha+1}}{(4\pi)^{\frac d2}}\,\frac d2\,\frac{\Gamma(\alpha-\frac d2-1)}{\Gamma(\alpha)}\,.
\end{split}
\end{equation}
In deriving \eqref{eq:dimregk2} the defining property of the $\Gamma$ function
$\Gamma(z+1)=z\,\Gamma(z)$ was used.

\subsection*{Renormalizations in \texorpdfstring{$\varphi^4$}{phi**4} theory, revisited}

Let us recall the basic idea of the (perturbative) renormalization program. To
develop renormalized perturbation theory, one starts with the action
\[
\mathcal A=\int d^dx\,\left[\frac12\,(\partial\varphi)^2+\frac{m^2}2\,\varphi^2+\frac\lambda{4!}\,\varphi^4+\text{counterterms}\right]\,,
\]
expressed in terms of the renormalized quantities. I have written it in the
$d$-dimensional form having in mind dimensional regularization. This action
contains counterterms
\[
\text{counterterms}=\frac{\delta Z}2\,(\partial\varphi)^2+\frac{\delta m^2}2\,\varphi^2+\frac{\delta\lambda}{4!}\,\varphi^4\,.
\]
We are planning to treat the theory perturbatively, with
\[
\mathcal A_I=\int_x\left[\frac\lambda{4!}\,\varphi^4+\text{counterterms}\right]
\]
taken as the perturbation.

The counterterms bring in additional vertices, so the diagrams will contain the
following elements (momentum $\delta$-functions are implicit)
\[
\dia{\draw[prop] (-1.5,0) -- (1.5,0); \draw[arr] (-0.5,0.35) -- (0.5,0.35) node[above,pos=0.5] {\small$p$};}=\frac1{p^2+m^2}\,,\qquad
\dia{\pic{xvertex};}=-\lambda\,,\qquad
\dia{\draw[prop,densely dashed] (-1.6,0) -- (1.6,0); \node[otimesv] at (0,0) {};}=-\bigl(p^2\,\delta Z+\delta m^2\bigr)\,,\qquad
\dia{\draw[prop] (-0.6,-0.6) -- (0.6,0.6); \draw[prop] (-0.6,0.6) -- (0.6,-0.6); \node[otimesv] at (0,0) {};}=-\delta\lambda\,.
\]
The counterterms $\delta Z$, $\delta m^2$, $\delta\lambda$ are assumed to be
power series in $\lambda$, i.e.
\begin{align*}
\delta Z&=z^{(1)}\,\lambda+z^{(2)}\,\lambda^2+\cdots\,,\\
\delta m^2&=b^{(1)}\,\lambda+b^{(2)}\,\lambda^2+\cdots\,,\\
\delta\lambda&=a^{(1)}\,\lambda+a^{(2)}\,\lambda^2+\cdots\,,
\end{align*}
with the coefficients dependent on $\Lambda$, or, in dimensional
regularization, on
\[
\epsilon=4-d\,,
\]
in such a way that all diagram contributions to the correlation functions of
renormalized fields have finite limit as $\Lambda\to\infty$, or, respectively,
$\epsilon\to0$ in dimensional regularization, order by order in $\lambda$. This
is general strategy of renormalized perturbation theory. Specific calculations
depend on the choice of the so called ``Renormalization Scheme''.

\subsection*{Renormalization Schemes}

Obviously, just satisfying this requirement (that the counterterms cancel all
singular contributions of ordinary diagrams) is not sufficient to fix the
counterterms uniquely. For instance, the renormalized field $\varphi$ always can
be renormalized further by multiplying it by a finite (i.e.\ regular at
$\epsilon=0$) constant
\[
\varphi=Z_{\text{finite}}^{\frac12}\,\tilde\varphi\,;
\]
similar finite renormalizations of $m^2$ and $\lambda$ are possible. Obviously,
such renormalizations do not spoil finiteness of renormalized perturbation
theory. We see that there is a whole class of renormalizable perturbation
theories, related one to another by finite renormalizations. These different
renormalized perturbation theories are called the \emph{Renormalization
Schemes}. Of course all these renormalized perturbation theories describe the
same field theory, and yield the correlation functions with identical physical
content. What is different, besides finite renormalizations of the correlation
functions due to the field renormalizations, is \emph{parametrization} of the
correlation functions. If $\varphi$, $m^2$, $\lambda$ and $\tilde\varphi$,
$\tilde m^2$, $\tilde\lambda$ are renormalized quantities in two different
renormalization schemes, the corresponding correlation functions are related as
\[
\Gamma^{(n)}(p_i|m^2,\lambda)=\Bigl[Z_{\text{finite}}(m^2,\lambda)\Bigr]^{-\frac n2}\,
\tilde\Gamma^{(n)}\bigl(p_i|\tilde m^2(m^2,\lambda),\tilde\lambda(m^2,\lambda)\bigr)\,.
\]
One can think of the parameters $(m^2,\lambda)$ as coordinates in the ``space
of $\varphi^4$ theories''. Different renormalization schemes correspond to
different coordinate systems in that space.

A renormalization scheme is usually defined by certain \emph{Normalization
conditions}.

\subsection*{Normalization conditions}

These are conditions that specify the choice of the renormalized parameters
$m^2$, $\lambda$, as well as the overall normalization of $\varphi$. This is
usually done by relating the renormalized parameters to the physical quantities
(this is not so in the so called ``minimal subtraction scheme'' which has more
technical definition; we will discuss it in some details later). In practice,
the schemes where these relations are simple are more convenient.

One such normalization condition we have discussed already when we defined the
renormalized mass parameter $m^2$ as the physical mass squared. Namely,
assuming that $\tilde\Gamma^{(2)}(p^2)$ has zero at some negative value of
$p^2$ we have defined $-m^2$ as the position of this zero. Note that this
condition does not depend on the overall normalization of $\varphi$. The latter
must be fixed by a separate condition. One can fix for instance the slope of
$\tilde\Gamma^{(2)}(p^2)$ at this zero $p^2=-m^2$; a convenient normalization
condition is to set this slope equal 1, i.e.
\begin{equation}\label{eq:SIa}
\tilde\Gamma^{(2)}(p^2)=p^2+m^2+O\bigl((p^2+m^2)^2\bigr)\qquad\text{as}\qquad p^2+m^2\to0\,.
\end{equation}
Equivalently, this condition means that the two-point correlation function
\begin{equation}\label{eq:SIW}
\tilde W(p)=\frac1{p^2+m^2}+O(1)\qquad\text{as}\qquad p^2+m^2\to0
\end{equation}
has a pole at $p^2=-m^2$ with the residue one. We also need to fix
normalization of $\lambda$. This is usually done by imposing certain condition
on $\tilde\Gamma^{(4)}$. The condition often taken in high-energy theory
calculations is
\begin{equation}\label{eq:SIb}
\tilde\Gamma^{(4)}(p,p,-p,-p)\Big|_{p^2=-m^2}=\lambda
\end{equation}
(the quantity in the left-hand side is related to $2\to2$ forward scattering
amplitude at zero energy). The equations \eqref{eq:SIW} and \eqref{eq:SIb} above
specify one possible renormalization scheme (I will refer to it as the scheme
SI).

Of course there are other schemes. Another scheme which is more frequently used
in statistical mechanics context is specified by the normalization conditions
\begin{align}
\tilde\Gamma^{(2)}(p^2)&=\tilde m^2+p^2+O(p^4)\,,\label{eq:SIIa0}\\
\tilde\Gamma^{(4)}(0,0,0,0)&=\tilde\lambda\,.\label{eq:SIIb0}
\end{align}
Here I have used the notations $\tilde m^2$ and $\tilde\lambda$ for the
associated renormalized parameters to stress that these parameters are
different from those defined in the scheme (SI) (in particular $\tilde m^2$
does not coincide with the physical mass), i.e.\ it defines different
``coordinates''. Henceforth this scheme is referred to as SII.

Let us see how this program works in few leading orders in the coupling
constant. The following calculation is equivalent to the one in PS, \S10.2,
10.5. There it is done using the normalization condition (SI). For the sake of
diversity I will instead adopt (SII). To facilitate the comparison I will use
notations $\tilde m^2$ and $\tilde\lambda$ for the renormalized parameters.
Also, to simplify notations I will write $\Gamma^{(n)}$ for the \emph{momentum
space} proper vertices (omitting the tilde).

\subsection*{Renormalization of \texorpdfstring{$\varphi^4$}{phi**4} at the leading (one loop) order}

Let me describe in some details the renormalization of the $\varphi^4$ theory
using dimensional regularization. For this illustration I will adopt the
renormalization scheme SII, as was described above. For reference, that scheme
is defined by the normalization conditions
\begin{align}
\tilde\Gamma^{(2)}(p^2)&=\tilde m^2+p^2+O(p^4)\,,\label{eq:SIIa}\\
\tilde\Gamma^{(4)}(0,0,0,0)&=\tilde\lambda\,.\label{eq:SIIb}
\end{align}
The normalization conditions specify the meaning of the renormalized parameters
$\tilde m^2$, $\tilde\lambda$, as well as the normalization of the field
$\varphi(x)$.

Consider first the two-point function,
\[
\Gamma^{(2)}(p^2)=p^2+\tilde m^2+\Sigma(p^2)\,.
\]
The normalization condition \eqref{eq:SIIa} states
\begin{equation}\label{eq:SIIsigma}
\Sigma(0)=0\,;\qquad\Sigma'(0)=0\,.
\end{equation}
In the leading order in $\tilde\lambda$ there are two diagrams contributing to
$\Sigma$,
\[
-\Bigl(\underbrace{\dia[0.4]{\draw[prop,densely dashed] (-2,0) -- (2,0); \draw[prop] (0,0.6) circle (0.6); \node[vtx] at (0,0) {};}}_{(a)}
+\underbrace{\ctline{1}}_{(b)}\Bigr)\,,
\]
where the diagram (b) comes from the $\sim\tilde\lambda$ terms in $\delta Z$ and
$\delta\tilde m^2$,
\[
\ctline{1}=-\tilde\lambda\,\bigl(z^{(1)}\,p^2+b^{(1)}\bigr)\,;
\]
The coefficients $z^{(1)}$ and $b^{(1)}$ are to be adjusted to satisfy
\eqref{eq:SIIsigma}.

We already have considered the diagram (a). It is a constant independent of
$p^2$; its value is
$(a)=-\frac{\tilde\lambda}2\int\frac{d^dk}{(2\pi)^d}\frac1{k^2+\tilde m^2}$.
Therefore to satisfy \eqref{eq:SIIsigma} we have to choose
\[
z^{(1)}=0\,;\qquad b^{(1)}=\frac{(a)}{\tilde\lambda}=-\frac12\int\frac{d^dk}{(2\pi)^d}\,\frac1{k^2+\tilde m^2}
=-\frac12\,\frac{(\tilde m^2)^{\frac d2-1}}{(4\pi)^{\frac d2}}\,\Gamma(1-d/2)\,,
\]
where the dimensional regularization was used to evaluate the integral. Note
that at this order there is a trivial cancellation between the diagrams (a) and
(b). Because of this, in considering higher-order diagrams we can in fact set
$b^{(1)}=0$ if we at the same time decide to ignore the ``bubble'' contributions
(a) in all propagators.

Consider now the leading perturbative corrections to the four-point function
$\Gamma^{(4)}$. At the tree level $\Gamma^{(4)}=\tilde\lambda$ and the
normalization condition \eqref{eq:SIIb} is satisfied at this order if we set
$a^{(1)}=0$, i.e.\ $\delta\tilde\lambda=O(\tilde\lambda^2)$. Nontrivial
contributions appear at the order $\tilde\lambda^2$; the corresponding diagrams
are
\[
\Gamma^{(4)}(p_1,\dots,p_4)=\tilde\lambda-\Bigl(\underset{(a)}{\bubS}+\underset{(b)}{\bubT}+\underset{(c)}{\bubU}+\underset{(d)}{\ctX{2}}\Bigr)
\]
where again the diagram (d) represents the contribution of the counterterm
$\delta\lambda$ in the order $\tilde\lambda^2$,
\begin{equation}\label{eq:a2vertex}
\ctX{2}=-\tilde\lambda^2\,a^{(2)}\,.
\end{equation}
Each of the diagrams (a-c) can be expressed through the integral
\begin{equation}\label{eq:Ip2}
I(p^2)=\int\frac{d^dk}{(2\pi)^d}\,\frac1{(k^2+\tilde m^2)\bigl((p+k)^2+\tilde m^2\bigr)}\,,
\end{equation}
namely
\[
(a)+(b)+(c)=\frac{\tilde\lambda^2}2\,\Bigl(I(p_{12}^2)+I(p_{13}^2)+I(p_{14}^2)\Bigr)\,,
\]
where $p_{12}=p_1+p_2$, $p_{13}=p_1+p_3$, $p_{14}=p_1+p_4$.

The normalization condition \eqref{eq:SIIb} requires that
\[
\Bigl[(a)+(b)+(c)+(d)\Bigr]_{p_i=0}=0\,,
\]
which fixes the value of the counterterm constant in \eqref{eq:a2vertex},
\[
a^{(2)}=\frac32\,I(0)\,.
\]
With this we obtain
\begin{equation}\label{eq:G4oneloop}
\Gamma^{(4)}(p_1,\dots,p_4)=\tilde\lambda-\frac{\tilde\lambda^2}2\,\Bigl(I_r(p_{12}^2)+I_r(p_{13}^2)+I_r(p_{14}^2)\Bigr)+O(\tilde\lambda^3)\,,
\end{equation}
where
\[
I_r(p^2)=I(p^2)-I(0)\,.
\]
Our goal now is to show that while the integral \eqref{eq:Ip2} is singular at
$d=4$ the combination appearing in \eqref{eq:G4oneloop} has finite limit at
$d=4$.

The integrand in \eqref{eq:Ip2} can be transformed using the so called
\emph{Feynman parametrization}. The identity\footnote{More general identity
\begin{equation}\label{eq:FeynParGen}
\frac1{A^\alpha}\,\frac1{B^\beta}=\frac{\Gamma(\alpha+\beta)}{\Gamma(\alpha)\Gamma(\beta)}\int_0^1\frac{dw\,w^{\alpha-1}(1-w)^{\beta-1}}{\bigl[wA+(1-w)B\bigr]^{\alpha+\beta}}
\end{equation}
is often useful. To verify it, change the integration variable to
$w=Bz/(Bz+A(1-z))$ to reduce the integral to the Euler's ``beta integral''
\[
B(\alpha,\beta):=\int_0^1dz\,z^{\alpha-1}(1-z)^{\beta-1}=\frac{\Gamma(\alpha)\Gamma(\beta)}{\Gamma(\alpha+\beta)}\,.
\]}
\[
\frac1{AB}=\int_0^1\frac{du}{\bigl(uA+(1-u)B\bigr)^2}
\]
which can be checked by explicit evaluation of the integral in the r.h.s.,
allows one to write the integral \eqref{eq:Ip2} as
\begin{equation}\label{eq:IFeyn}
I(p^2)=\int_0^1du\int\frac{d^dk}{(2\pi)^d}\,\frac1{\bigl(\tilde m^2+k^2+2u\,(kp)+u\,p^2\bigr)^2}\,.
\end{equation}
The combination in the denominator can be put in the form
\[
\tilde m^2+k^2+2u\,(kp)+u\,p^2=\tilde m^2+u(1-u)\,p^2+(k+up)^2\,,
\]
and by applying the translation property (i) of the dimensionally continued
integrals we have
\[
I(p^2)=\int_0^1du\int\frac{d^dk}{(2\pi)^d}\,\frac1{\bigl(k^2+u(1-u)\,p^2+\tilde m^2\bigr)^2}\,.
\]
The momentum integral has the form \eqref{eq:dimregalpha} with $\alpha=2$, hence
\[
I(p^2)=\frac{\Gamma(2-\frac d2)}{(4\pi)^{\frac d2}}\int_0^1du\,\bigl[\tilde m^2+u(1-u)\,p^2\bigr]^{\frac d2-2}\,.
\]
We see that due to the factor $\Gamma(2-d/2)=2/\epsilon-\gamma+O(\epsilon)$,
$\epsilon:=4-d$, this expression has pole at $d=4$. However the combination
\[
I_r(p^2)=I(p^2)-I(0)=\frac{\Gamma(2-\frac d2)}{(4\pi)^{\frac d2}}\int_0^1du\,\Bigl(\bigl[\tilde m^2+u(1-u)\,p^2\bigr]^{\frac d2-2}-\bigl[\tilde m^2\bigr]^{\frac d2-2}\Bigr)
\]
has finite value at $d=4$ as the result of compensation between the pole of the
gamma function and zero of the integrand. One finds
\begin{align*}
I_r(p^2)\Big|_{d=4}&=-\frac1{(4\pi)^2}\int_0^1du\,\log\left(1+u(1-u)\,\frac{p^2}{\tilde m^2}\right)\\
&=-\frac1{16\pi^2}\left[\sqrt{\frac{p^2+4\tilde m^2}{p^2}}\,\log\frac{\sqrt{p^2+4\tilde m^2}+\sqrt{p^2}}{\sqrt{p^2+4\tilde m^2}-\sqrt{p^2}}-2\right]\,.
\end{align*}
I leave evaluating the last explicit expression as an \textbf{Exercise}. We see
that by giving the counterterm coefficient $a^{(2)}$ certain dependence on $d$
which is singular at $d=4$,
\begin{equation}\label{eq:a2}
a^{(2)}=\frac32\,I(0)=\frac32\,\frac{\Gamma(2-\frac d2)}{(4\pi)^{\frac d2}}\,(\tilde m^2)^{\frac d2-2}\,,
\end{equation}
we compensated the singularities of the diagrams (a), (b), (c) above.

\begin{exercise}
Compare the above calculation with the one in \S10.2 of PS. Find the relation
(to the order $\lambda^2$) between our renormalized parameters $\tilde m^2$,
$\tilde\lambda$ and parameters $m^2$, $\lambda$ defined there, which correspond
to the scheme (SI).
\end{exercise}

It is instructive to study renormalization at the next, two loop order. The two
loop contributions correspond to the terms $\sim\lambda^2$ in $\Gamma^{(2)}$
and $\lambda^3$ in $\Gamma^{(4)}$. Analysis of the two loop terms in
$\Gamma^{(4)}$ can be found in PS Sect.~10.5. I repeat it in the Appendix
below, along with the brief discussion of the two-loop divergences in
$\Gamma^{(2)}$.

To complete our analysis, we have to check that no new counterterms are
required to compensate for divergences in higher proper vertices
$\Gamma^{(6)}$, $\Gamma^{(8)}$, etc. At one-loop level the higher vertices are
finite without any counterterms since the diagrams
\[
-\Gamma^{(2n)}_{\text{1-loop}}=\diac[0.62]{\draw[prop] (0,0) circle (1.1);
  \foreach \a in {20,110,200,290} {\node[vtx] at (\a:1.1) {}; \draw[prop,densely dashed] (\a:1.1) -- ++(\a+30:0.9); \draw[prop,densely dashed] (\a:1.1) -- ++(\a-30:0.9);}
  \node[font=\tiny] at (245:1.6) {$\cdot\,\cdot$};}
\]
are convergent at $d=4$. The two-loop contributions to these vertices involve
divergent diagrams, with the divergences coming from the 2- or 4-leg
subdiagrams. For example, $\Gamma^{(6)}_{\text{2-loop}}$ receives contribution
from the diagram
\[
-\Gamma^{(6)}=\cdots+\diac[0.62]{\draw[prop] (0,0) circle (1.1);
  \draw[prop] (125:1.1) .. controls (0,0.2) .. (55:1.1);
  \foreach \a/\b in {125/150,55/30} {\node[vtx] at (\a:1.1) {}; \draw[prop,densely dashed] (\a:1.1) -- ++(\b:0.9);}
  \foreach \a/\b/\c in {200/180/240,320/290/350} {\node[vtx] at (\a:1.1) {}; \draw[prop,densely dashed] (\a:1.1) -- ++(\b:0.9); \draw[prop,densely dashed] (\a:1.1) -- ++(\c:0.9);}}
\]
which is singular as $1/\epsilon$. However, it is possible to check that the
one-loop counterterms
\[
\frac{\lambda\,b^{(1)}(\epsilon)}2\,\varphi^2+\frac{\lambda^2\,a^{(2)}(\epsilon)}{4!}\,\varphi^4
\]
are sufficient to cancel all such divergences. For instance, the above divergent
diagram comes along with two other divergent diagrams
\[
\diac[0.62]{\draw[prop] (0,0) circle (1.1); \draw[prop] (125:1.1) .. controls (0,0.2) .. (55:1.1);
  \foreach \a/\b/\l in {125/150/1,55/30/2} {\node[vtx] at (\a:1.1) {}; \draw[prop,densely dashed] (\a:1.1) -- ++(\b:0.9) node[pos=1.3,font=\scriptsize] {\l};}
  \foreach \a/\b/\c/\l/\m in {200/180/240/6/5,320/290/350/4/3} {\node[vtx] at (\a:1.1) {}; \draw[prop,densely dashed] (\a:1.1) -- ++(\b:0.9) node[pos=1.3,font=\scriptsize] {\l}; \draw[prop,densely dashed] (\a:1.1) -- ++(\c:0.9) node[pos=1.3,font=\scriptsize] {\m};}}
\qquad
\diac[0.62]{\draw[prop] (0,0) circle (1.1); \draw[prop] (125:1.1) .. controls (0,0.2) .. (55:1.1);
  \foreach \a/\b/\l in {125/150/2,55/30/1} {\node[vtx] at (\a:1.1) {}; \draw[prop,densely dashed] (\a:1.1) -- ++(\b:0.9) node[pos=1.3,font=\scriptsize] {\l};}
  \foreach \a/\b/\c/\l/\m in {200/180/240/6/5,320/290/350/4/3} {\node[vtx] at (\a:1.1) {}; \draw[prop,densely dashed] (\a:1.1) -- ++(\b:0.9) node[pos=1.3,font=\scriptsize] {\l}; \draw[prop,densely dashed] (\a:1.1) -- ++(\c:0.9) node[pos=1.3,font=\scriptsize] {\m};}}
\qquad
\diac[0.62]{\draw[prop] (0,0) circle (1.1); \draw[prop] (0,1.55) circle (0.45);
  \node[vtx] at (0,1.1) {}; \node[vtx] at (0,2.0) {};
  \draw[prop,densely dashed] (0,2.0) -- ++(135:0.9) node[pos=1.3,font=\scriptsize] {1}; \draw[prop,densely dashed] (0,2.0) -- ++(45:0.9) node[pos=1.3,font=\scriptsize] {2};
  \foreach \a/\b/\c/\l/\m in {200/180/240/6/5,320/290/350/4/3} {\node[vtx] at (\a:1.1) {}; \draw[prop,densely dashed] (\a:1.1) -- ++(\b:0.9) node[pos=1.3,font=\scriptsize] {\l}; \draw[prop,densely dashed] (\a:1.1) -- ++(\c:0.9) node[pos=1.3,font=\scriptsize] {\m};}}
\]
The sum of the singular part of these three diagrams is canceled by the diagram
\[
\diac[0.62]{\draw[prop] (0,0) circle (1.1);
  \node[otimesv] at (0,1.1) {};
  \draw[prop,densely dashed] (0,1.1) -- ++(135:0.9) node[pos=1.3,font=\scriptsize] {1}; \draw[prop,densely dashed] (0,1.1) -- ++(45:0.9) node[pos=1.3,font=\scriptsize] {2};
  \foreach \a/\b/\c/\l/\m in {200/180/240/6/5,320/290/350/4/3} {\node[vtx] at (\a:1.1) {}; \draw[prop,densely dashed] (\a:1.1) -- ++(\b:0.9) node[pos=1.3,font=\scriptsize] {\l}; \draw[prop,densely dashed] (\a:1.1) -- ++(\c:0.9) node[pos=1.3,font=\scriptsize] {\m};}}
\]
involving the one-loop counterterm vertex.

Similar situation persists to higher loops. Suppose we have renormalized the
theory (i.e.\ determined all the counterterms as the functions of $\epsilon$) up
to $L-1$ loops. Then contributions of all $L$-loop diagrams, with all
$L-1$-loop counterterms, to the effective action $\Gamma[\phi]$ can be written
as
\[
\Gamma^{\text{div}}_L[\phi]+\Gamma^{\text{finite}}_L[\phi]\,,
\]
where $\Gamma^{\text{div}}_L$ has the form
\[
-\tilde\lambda^L\int d^dx\,\left[\frac{Z^{(L)}(\epsilon)}2\,(\partial\phi)^2+\frac{B^{(L)}(\epsilon)}2\,\phi^2+\frac{\tilde\lambda\,A^{(L)}(\epsilon)}{4!}\,\phi^4\right]\,,
\]
with coefficients having $L$th order pole in $\epsilon$,
\begin{equation}\label{eq:divpoles}
Z^{(L)}=\sum_{k=1}^{L-1}\frac{Z^{(L)}_{-k}}{\epsilon^k}\,;\qquad
B^{(L)}=\sum_{k=1}^L\frac{B^{(L)}_{-k}}{\epsilon^k}\,;\qquad
A^{(L)}=\sum_{k=1}^L\frac{A^{(L)}_{-k}}{\epsilon^k}\,,
\end{equation}
and $\Gamma^{\text{finite}}$ is finite at $\epsilon=0$. It is clear then that the
divergent part can be compensated for if one adds to the action the $L$-loop
counterterm
\[
\mathcal A^{(L)}=\tilde\lambda^L\int d^dx\,\left[\frac{z^{(L)}(\epsilon)}2\,(\partial\varphi)^2+\frac{b^{(L)}(\epsilon)}2\,\varphi^2+\frac{\tilde\lambda\,a^{(L)}(\epsilon)}{4!}\,\varphi^4\right]\,,
\]
with
\[
z^{(L)}(\epsilon)=Z^{(L)}(\epsilon)+Z_0^{(L)}\,;\qquad
b^{(L)}(\epsilon)=B^{(L)}(\epsilon)+B_0^{(L)}\,;\qquad
a^{(L)}(\epsilon)=A^{(L)}(\epsilon)+A_0^{(L)}\,,
\]
where $Z_0^{(L)},B_0^{(L)},A_0^{(L)}$ are finite numbers.

The freedom in the choice of these finite numbers represents the freedom in the
choice of the renormalization scheme. In given scheme they are determined by the
normalization conditions, like \eqref{eq:SIW}, \eqref{eq:SIb} or
\eqref{eq:SIIa0}, \eqref{eq:SIIb0}, which usually directly relate the parameters
$\lambda$ and $m^2$ to physical quantities. The form \eqref{eq:divpoles} of the
divergent part of the effective action suggests another class of renormalization
schemes. In so called \emph{Minimal Subtraction Scheme} one simply sets
$Z_0^{(L)}=0$, $B_0^{(L)}=0$, $A_0^{(L)}=0$. In other words, in this scheme one
chooses the $L$-loop counterterm to be exactly equal to the divergent part (i.e.\
the singular part of the Laurent expansion around $\epsilon=0$) of the effective
action renormalized to $L-1$ loops, as defined above, i.e.
\[
\mathcal A^{(L)}[\varphi]=\Gamma^{\text{div}}_L[\varphi]\,.
\]
In this scheme the parameters $m^2_{MS}$ and $\lambda_{MS}$ entering the action
\[
\mathcal A=\int d^dx\,\left[\frac12\,(\partial\varphi)^2+\frac{m^2_{MS}}2\,\varphi^2+\frac{\lambda_{MS}}{4!}\,\varphi^4+\text{Counterterms in MS Scheme}\right]\,,
\]
are not a priori related to physical quantities such as the particle mass and
scattering amplitudes. This relation has to be calculated order by order in
perturbation theory.

Full proof that such procedure works to all orders is rather involved (it can be
found e.g.\ in \emph{J.~Zinn-Justin, Quantum Field Theory and Critical
Phenomena}, or in original papers quoted on p.~338 of PS). According to these
general theorems similar renormalization procedure works to all orders in any
field theory whose perturbation theory has only finitely many primitive
divergences. We have seen that this is true if and only if the mass dimensions
of all couplings are non-negative.

\subsection*{Appendix A: Two loop analysis of \texorpdfstring{$\varphi^4$}{phi**4} (PS Sect.~10.5)}

Here we carry out two-loop renormalization analysis in the $\varphi^4$ theory.
Let us now start with the $\tilde\lambda^3$ contributions to $\Gamma^{(4)}$. I
will still assume the renormalization scheme SII. Also, I will disregard the
$b^{(1)}$ term in $\delta m^2$ counterterm and the ``trivial bubble diagram''; we
saw that these cancel each other exactly. Then we have 16 diagrams contributing
to $\Gamma^{(4)}$ in this order:
\begin{multline*}
-\Gamma^{(4)}(p_1,\dots,p_4)=\cdots+\Bigl(\underset{(a)}{\diac[0.4]{\coordinate (L) at (-1.6,0); \coordinate (M) at (0,0); \coordinate (R) at (1.6,0);
  \draw[prop] (-0.8,0) ellipse (0.8 and 0.5); \draw[prop] (0.8,0) ellipse (0.8 and 0.5);
  \legs{L/135/1,L/225/2,R/45/3,R/315/4} \foreach \p in {L,M,R} \node[vtx] at (\p) {};}}
+\underset{(b)}{\diac[0.4]{\coordinate (L) at (-0.8,0); \coordinate (T) at (0.55,0.6); \coordinate (B) at (0.55,-0.6);
  \draw[prop] (0,0) circle (0.8); \draw[prop] (T) .. controls (0.25,0) .. (B);
  \legs{L/135/1,L/225/2,T/45/3,B/315/4} \foreach \p in {L,T,B} \node[vtx] at (\p) {};}}
+\underset{(c)}{\diac[0.4]{\coordinate (R) at (0.8,0); \coordinate (T) at (-0.55,0.6); \coordinate (B) at (-0.55,-0.6);
  \draw[prop] (0,0) circle (0.8); \draw[prop] (T) .. controls (-0.25,0) .. (B);
  \legs{T/135/1,B/225/2,R/45/3,R/315/4} \foreach \p in {R,T,B} \node[vtx] at (\p) {};}}\\
+\underset{(d)}{\diac[0.4]{\coordinate (L) at (-0.9,0); \coordinate (R) at (0.9,0);
  \draw[prop] (0,0) ellipse (0.9 and 0.55); \legs{L/135/1,L/225/2,R/45/3,R/315/4}
  \node[vtx] at (L) {}; \node[otimesv,minimum size=5pt] at (R) {}; \node[font=\tiny] at (1.25,0) {(2)};}}
+\underset{(e)}{\diac[0.4]{\coordinate (L) at (-0.9,0); \coordinate (R) at (0.9,0);
  \draw[prop] (0,0) ellipse (0.9 and 0.55); \legs{L/135/1,L/225/2,R/45/3,R/315/4}
  \node[otimesv,minimum size=5pt] at (L) {}; \node[vtx] at (R) {}; \node[font=\tiny] at (-1.25,0) {(2)};}}\Bigr)\\
+\Bigl(\text{10 diagrams with permutations of the external legs}\Bigr)+\Bigl(\ctX{3}\Bigr)+\cdots
\end{multline*}
The diagrams (d) and (e) (and 4 similar diagrams with permutations of the
external legs) involve the counterterm vertex $-\tilde\lambda^2a^{(2)}$ with the
coefficient $a^{(2)}$ which was already determined in the one-loop analysis, see
Eq.~\eqref{eq:a2} above.

The last diagram shown above represents the contribution of the term
$\tilde\lambda^3$ in $\delta\lambda$ counterterm. We must see that it is indeed
possible to adjust this single coefficient $a^{(3)}$ in such a way that all
divergences (i.e.\ all singularities at $d=4$) cancel. This last contribution is
momentum independent constant. For our renormalization program to be successful
to this order we must prove that all momentum-dependent divergences cancel out
among the first 15 diagrams above.

Consider 5 diagrams (a-e) above. It is useful to divide them into three groups
as follows (the counterterm diagrams (d) and (e) are shared among the groups,
with the weights shown)
\begin{align*}
\text{I)}&\quad\diac[0.4]{\coordinate (L) at (-1.6,0); \coordinate (M) at (0,0); \coordinate (R) at (1.6,0);
  \draw[prop] (-0.8,0) ellipse (0.8 and 0.5); \draw[prop] (0.8,0) ellipse (0.8 and 0.5);
  \legs{L/135/1,L/225/2,R/45/3,R/315/4} \foreach \p in {L,M,R} \node[vtx] at (\p) {};}
+\frac13\,\diac[0.4]{\coordinate (L) at (-0.9,0); \coordinate (R) at (0.9,0);
  \draw[prop] (0,0) ellipse (0.9 and 0.55); \legs{L/135/1,L/225/2,R/45/3,R/315/4}
  \node[otimesv,minimum size=5pt] at (L) {}; \node[vtx] at (R) {}; \node[font=\tiny] at (-1.25,0) {(2)};}
+\frac13\,\diac[0.4]{\coordinate (L) at (-0.9,0); \coordinate (R) at (0.9,0);
  \draw[prop] (0,0) ellipse (0.9 and 0.55); \legs{L/135/1,L/225/2,R/45/3,R/315/4}
  \node[vtx] at (L) {}; \node[otimesv,minimum size=5pt] at (R) {}; \node[font=\tiny] at (1.25,0) {(2)};}\,,\\
\text{II)}&\quad\diac[0.4]{\coordinate (R) at (0.8,0); \coordinate (T) at (-0.55,0.6); \coordinate (B) at (-0.55,-0.6);
  \draw[prop] (0,0) circle (0.8); \draw[prop] (T) .. controls (-0.25,0) .. (B);
  \legs{T/135/1,B/225/2,R/45/3,R/315/4} \foreach \p in {R,T,B} \node[vtx] at (\p) {};}
+\frac23\,\diac[0.4]{\coordinate (L) at (-0.9,0); \coordinate (R) at (0.9,0);
  \draw[prop] (0,0) ellipse (0.9 and 0.55); \legs{L/135/1,L/225/2,R/45/3,R/315/4}
  \node[otimesv,minimum size=5pt] at (L) {}; \node[vtx] at (R) {}; \node[font=\tiny] at (-1.25,0) {(2)};}\,,\\
\text{III)}&\quad\diac[0.4]{\coordinate (L) at (-0.8,0); \coordinate (T) at (0.55,0.6); \coordinate (B) at (0.55,-0.6);
  \draw[prop] (0,0) circle (0.8); \draw[prop] (T) .. controls (0.25,0) .. (B);
  \legs{L/135/1,L/225/2,T/45/3,B/315/4} \foreach \p in {L,T,B} \node[vtx] at (\p) {};}
+\frac23\,\diac[0.4]{\coordinate (L) at (-0.9,0); \coordinate (R) at (0.9,0);
  \draw[prop] (0,0) ellipse (0.9 and 0.55); \legs{L/135/1,L/225/2,R/45/3,R/315/4}
  \node[vtx] at (L) {}; \node[otimesv,minimum size=5pt] at (R) {}; \node[font=\tiny] at (1.25,0) {(2)};}\,.
\end{align*}
Let us show that all momentum dependent divergences cancel separately within
each group.

Consider first the group I. The first diagram there is easily calculated to be
\[
-\frac{\tilde\lambda^3}4\,\bigl[I(p_{12}^2)\bigr]^2\,,
\]
the 4 in the denominator being the symmetry factor, $p_{12}=p_1+p_2$, and again
\[
I(p^2)=\int\frac{d^dk}{(2\pi)^d}\,\frac1{(k^2+\tilde m^2)\bigl((p+k)^2+\tilde m^2\bigr)}\,.
\]
Now, each of the other two diagrams in this group contribute
\[
\frac13\,\frac{\tilde\lambda^3}2\,a^{(2)}\,I(p_{12}^2)\,,
\]
so that the sum of all three diagrams in this group is
\[
\text{Group I}=-\frac{\tilde\lambda^3}4\,\bigl[I(p_{12}^2)\bigr]^2+\frac23\,\frac{\tilde\lambda^3}2\,\frac32\,I(0)\,I(p_{12}^2)
=-\frac{\tilde\lambda^3}4\,\bigl[I(p_{12}^2)-I(0)\bigr]^2+\frac{\tilde\lambda^3}4\,\bigl[I(0)\bigr]^2
\]
where the expression $a^{(2)}=\frac32I(0)$, Eq.~\eqref{eq:a2}, is used. We know
from the one-loop analysis that the difference $I_r(p^2)=I(p^2)-I(0)$ is finite
at $d=4$ and therefore the divergent parts of the diagrams in the group I indeed
sum up to a momentum independent constant (the last term in the expression
above). Note that this divergent constant has second order pole at $d=4$, i.e.\
it behaves as $1/\epsilon^2$ as $\epsilon\to0$.

Consider now the group II. It contains the diagram
\[
\diac[0.45]{\coordinate (R) at (0.8,0); \coordinate (T) at (-0.55,0.6); \coordinate (B) at (-0.55,-0.6);
  \draw[prop] (0,0) circle (0.8); \draw[prop] (T) .. controls (-0.25,0) .. (B);
  \foreach \p/\a/\n in {T/135/p_2,B/225/p_1,R/45/p_3,R/315/p_4} {\draw[prop,densely dashed] (\p) -- ++(\a:1.1) node[pos=1.25,font=\tiny] {$\n$};}
  \draw[arr] (60:1.15) arc[start angle=60,end angle=100,radius=1.15]; \node[font=\tiny] at (80:1.45) {$k$};
  \draw[arr] (260:1.15) arc[start angle=260,end angle=300,radius=1.15]; \node[font=\tiny] at (280:1.5) {$p_{12}+k$};
  \foreach \p in {R,T,B} \node[vtx] at (\p) {};}
=-\frac{\tilde\lambda^3}2\int\frac{d^dk}{(2\pi)^d}\,\frac1{k^2+\tilde m^2}\,\frac1{(p_{12}+k)^2+\tilde m^2}\,I\bigl((k+p_1)^2\bigr)\,,
\]
where $I(p^2)$ is the integral \eqref{eq:Ip2}, which we already found equal to
\[
I(p^2)=\frac{\Gamma(2-d/2)}{(4\pi)^{\frac d2}}\int_0^1du\,\Bigl[\tilde m^2+u(1-u)\,p^2\Bigr]^{\frac d2-2}\,.
\]
Using this formula and combining the propagators with the help of Feynman
parametrization we obtain for the above diagram
\begin{multline*}
-\frac{\tilde\lambda^3}2\,\frac{\Gamma(2-d/2)}{(4\pi)^{\frac d2}}
\int_0^1du\int_0^1dv\int\frac{d^dk}{(2\pi)^d}\,\frac1{\bigl(k^2+2v\,(kp_{12})+v\,p_{12}^2+\tilde m^2\bigr)^2}\\
\times\frac1{\bigl(\tilde m^2+u(1-u)\,(k+p_1)^2\bigr)^{2-\frac d2}}\,.
\end{multline*}
The $k$ integral here can be evaluated by using the transformation generalizing
the Feynman parametrization. Namely, consider the identity
\eqref{eq:FeynParGen},
\begin{equation}\label{eq:FeynParGen2}
\frac1{A^\alpha\,B^\beta}=\frac{\Gamma(\alpha+\beta)}{\Gamma(\alpha)\Gamma(\beta)}\int_0^1\frac{dw\,w^{\alpha-1}(1-w)^{\beta-1}}{\bigl[wA+(1-w)B\bigr]^{\alpha+\beta}}\,.
\end{equation}
It can be derived by explicit evaluation of the integral in the right-hand side:
projective transformation of the integration variable $w=Bz/(Bz+A(1-z))$ brings
it to the form of the Beta-integral
\[
B(\alpha,\beta)=\int_0^1dz\,z^{\alpha-1}(1-z)^{\beta-1}=\frac{\Gamma(\alpha)\Gamma(\beta)}{\Gamma(\alpha+\beta)}
\]
(details can be found in \S\,10.5 of PS). Using this identity (with
$\alpha=2-\frac d2$, $\beta=2$) one can write the above integral as
\begin{multline*}
-\frac{\tilde\lambda^3}2\,\frac{\Gamma(4-d/2)}{(4\pi)^{\frac d2}}\int_0^1du\,dv\,dw\\
\times\int\frac{d^dk}{(2\pi)^d}\,\frac{w^{1-\frac d2}\,(1-w)}{\Bigl[w\bigl(\tilde m^2+u(1-u)\,(k+p_1)^2\bigr)+(1-w)\bigl(\tilde m^2+k^2+2v\,(kp_{12})+v\,p_{12}^2\bigr)\Bigr]^{4-\frac d2}}\,.
\end{multline*}
By appropriate shift of the $d$-momentum integration variable the denominator
can be brought to the form
\[
A\,q^2+P^2+\tilde m^2
\]
where
\[
A=(1-w)+w\,u(1-u)\,,\qquad A\big|_{w=0}=1\,,
\]
and $P$ is complicated function of $p_{12}$, $p_1$ and the parameters $u,v,w$.
We only need to know that
\[
P^2\Big|_{w=0}=v(1-v)\,p_{12}^2\,.
\]
In this representation the integral over $q$ is readily evaluated,
\[
-\frac{\tilde\lambda^3}2\,\frac{\Gamma(4-d)}{(4\pi)^d}\int_0^1du\,dv\,dw\,\frac{w^{1-\frac d2}\,(1-w)}{A^{\frac d2}}\,\bigl(\tilde m^2+P^2\bigr)^{d-4}\,.
\]
Evidently this expression is singular at $d=4$. First, there is a pole at $d=4$
in the factor
\[
\Gamma(4-d)=\Gamma(\epsilon)=\frac1\epsilon+O(1)
\]
in front. In addition, at $d=4$ the integral over $w$ diverges at $w\to0$
because of the factor $w^{1-\frac d2}$ in the integrand; as the result the
integral over $w$ has additional pole at $d=4$, i.e.\ the whole diagram again
has a second order pole at this point.

To isolate the singular contributions let us write the integral over $w$ as
\[
\int_0^1dw\,w^{1-\frac d2}\,f(w)=f(0)\int_0^1dw\,w^{1-\frac d2}+\int_0^1dw\,w^{1-\frac d2}\,\bigl[f(w)-f(0)\bigr]\,.
\]
The second piece is
\[
-\frac{\tilde\lambda^3}2\,\frac{\Gamma(4-d)}{(4\pi)^d}\int_0^1du\,dv\,dw\,w^{1-\frac d2}\,
\left[\frac{1-w}{A^{\frac d2}}\,\bigl(\tilde m^2+P^2(w)\bigr)^{d-4}-\bigl(\tilde m^2+P^2(w=0)\bigr)^{d-4}\right]\,.
\]
Here the integral over $u,v,w$ is convergent and brings in no additional
singularity at $d=4$. Therefore the above expression has only simple pole at
$\epsilon=0$. To evaluate the residue at this simple pole one can set $d=4$
everywhere except in $\Gamma(4-d)$. It is clear then that this residue does not
depend on external momenta, i.e.\ this part of the above diagram has the form
\[
\frac{C_0}\epsilon+C_1(p_1,\dots,p_4)\,,
\]
where $C_0$ does not depend on $p_i$, and the $C_1$ is finite at $\epsilon=0$.

Let us take a look at the other piece, proportional to $f(0)$. It is
\[
f(0)\int_0^1dw\,w^{1-\frac d2}=\frac2\epsilon\,f(0)=-\frac2\epsilon\,\frac{\tilde\lambda^3}2\,\frac{\Gamma(\epsilon)}{(4\pi)^d}\int_0^1du\,dv\,\bigl(\tilde m^2+v(1-v)\,p_{12}^2\bigr)^{-\epsilon}
\]
The $u$ integral here is trivial. Now, for $\epsilon\to0$ the integral can be
written as
\[
\int_0^1dv\,\Bigl(1-\epsilon\,\log\bigl(\tilde m^2+v(1-v)\,p_{12}^2\bigr)+O(\epsilon^2)\Bigr)\,,
\]
and therefore this piece of the diagram is
\[
-\frac{\tilde\lambda^3}{(4\pi)^4}\,\frac1\epsilon\left(\frac1\epsilon-\gamma+\log4\pi-\int_0^1dv\,\log\bigl[\tilde m^2+v(1-v)\,p_{12}^2\bigr]\right)+\text{finite terms.}
\]
We have to add here the second diagram in the group II, which is
\[
\frac23\,\frac{\tilde\lambda^3}2\,a^{(2)}\,I(p_{12}^2)=\frac{\tilde\lambda^3}2\,I(0)\,I(p_{12}^2)\,,
\]
where the expression \eqref{eq:a2} for the coefficient $a^{(2)}$ is used. Recall
that
\begin{align*}
I(p^2)&=\frac{\Gamma(2-\frac d2)}{(4\pi)^{\frac d2}}\int_0^1du\,\bigl[\tilde m^2+u(1-u)\,p^2\bigr]^{\frac d2-2}\\
&=\frac1{(4\pi)^{\frac d2}}\left(\frac2\epsilon-\gamma-\int_0^1du\,\log\bigl(\tilde m^2+u(1-u)\,p^2\bigr)+O(\epsilon)\right)\,.
\end{align*}
Therefore the contribution of this diagram is
\[
\frac{\tilde\lambda^3}{(4\pi)^4}\,\frac1\epsilon\left(\frac2\epsilon-2\gamma+2\log4\pi-\log\tilde m^2-\int_0^1du\,\log\bigl[\tilde m^2+u(1-u)\,p_{12}^2\bigr]\right)+\text{finite terms.}
\]
Adding the two diagrams of the group II we see that the $p_{12}$-dependent
singular terms cancel out,
\[
\text{Group II}=\frac{\tilde\lambda^3}{(4\pi)^4}\,\frac1\epsilon\left(\frac1\epsilon-\gamma+\log4\pi-\log\tilde m^2\right)+\text{finite terms.}
\]
The group III can be analyzed similarly. We find that all singular parts of the
first 15 diagrams are pole terms with the coefficients independent of external
momenta,
\begin{equation}\label{eq:15diagrams}
\text{15 diagrams}=\tilde\lambda^3\left(\frac{c_{-2}}{\epsilon^2}+\frac{c_{-1}}\epsilon+F(p_1,\dots,p_4)\right)\,,
\end{equation}
where $c_{-2},c_{-1}$ do not depend on $p_i$, and $F$ is finite at $d=4$.
Therefore the counterterm coefficient $a^{(3)}$ appearing in the 16th diagram
always can be appropriately adjusted to compensate the singular terms. In fact,
according to our normalization condition \eqref{eq:SIIb}
($\Gamma^{(4)}(0,0,0,0)=\tilde\lambda$), this coefficient is determined as
\[
a^{(3)}=\frac{c_{-2}}{\epsilon^2}+\frac{c_{-1}}\epsilon+F(0,\dots,0)
\]
and hence
\[
\text{All 16 diagrams}=\tilde\lambda^3\Bigl(F(p_1,\dots,p_4)-F(0,\dots,0)\Bigr)\,.
\]

\begin{exercise}
Verify these two-loop calculations. Evaluate $c_{-2}$ and $c_{-1}$.
\end{exercise}

Let us briefly consider the $\tilde\lambda^2$ contributions to $\Gamma^{(2)}$.
These are given by three diagrams
\[
\Gamma^{(2)}(p^2)=p^2+\tilde m^2-\Bigl(\underset{(a)}{\diac[0.45]{\draw[prop,densely dashed] (-2.3,0) -- (-1.1,0) (1.1,0) -- (2.3,0);
  \draw[prop] (-1.1,0) -- (1.1,0); \draw[prop] (0,0) ellipse (1.1 and 0.7);
  \node[above,font=\tiny] at (0,0.75) {$p+k_1+k_2$}; \node[font=\tiny] at (0,0.25) {$k_1$}; \node[below,font=\tiny] at (0,-0.75) {$k_2$};
  \node[vtx] at (-1.1,0) {}; \node[vtx] at (1.1,0) {};}}
+\underset{(b)}{\diac[0.45]{\draw[prop,densely dashed] (-1.8,0) -- (1.8,0); \draw[prop] (0,0.6) circle (0.6);
  \node[otimesv,minimum size=5pt] at (0,0) {}; \node[font=\tiny] at (0,0.6) {(2)};}}
+\underset{(c)}{\ctline{2}}\Bigr)
\]
The diagram (b) involves the counterterm vertex $-\tilde\lambda^2a^{(2)}$, and
the diagram (c) contains the $\tilde\lambda^2$ part of the counterterms $\delta Z$
and $\delta m^2$,
\[
\ctline{2}=-\tilde\lambda^2\,\bigl(z^{(2)}\,p^2+b^{(2)}\bigr)\,.
\]
Again, the coefficients here must be determined from the normalization condition
\eqref{eq:SIIa}.

Let us denote $\Sigma_a(p^2)$ the contribution of the diagram (a) to the mass
operator. One can write it as
\[
\Sigma_a(p^2)=\Sigma_a(0)+p^2\,\Sigma_a'(0)+\Sigma_r(p^2)\,,
\]
where by the definition
\[
\Sigma_r(p^2)=O(p^4)\qquad\text{as}\qquad p^2\to0\,.
\]
Let us analyse the diagram (a) and show that $\Sigma_r(p^2)$ has finite value at
$d=4$. The contribution of the diagram (a) is given by the integral
\[
\Sigma_a(p^2)=-\frac{\tilde\lambda^2}{3!}\int\frac{d^dk_1\,d^dk_2}{(2\pi)^{2d}}\,\frac1{(k_1^2+\tilde m^2)(k_2^2+\tilde m^2)\bigl((p+k_1+k_2)^2+\tilde m^2\bigr)}\,.
\]
One of the $k$-integrations, say over $k_2$, can be done exactly the way we did
one-loop integrals for $\Gamma^{(4)}$:
\begin{align*}
&\int\frac{d^dk_2}{(2\pi)^d}\,\frac1{(k_2^2+\tilde m^2)\bigl((p+k_1+k_2)^2+\tilde m^2\bigr)}=I\bigl((p+k_1)^2\bigr)\\
&\qquad=\frac{\Gamma(2-d/2)}{(4\pi)^{\frac d2}}\int_0^1du\,\Bigl[\tilde m^2+u(1-u)\,(p+k_1)^2\Bigr]^{\frac d2-2}\,.
\end{align*}
We find
\[
\Sigma_a(p^2)=-\frac{\tilde\lambda^2}{3!}\,\frac{\Gamma(2-d/2)}{(4\pi)^{\frac d2}}\int_0^1du\int\frac{d^dk_1}{(2\pi)^d}\,
\frac1{\bigl(u(1-u)\,(p+k_1)^2+\tilde m^2\bigr)^{2-d/2}\,(k_1^2+\tilde m^2)}\,.
\]
The remaining integral over $k_1$ can be done using \eqref{eq:FeynParGen2} which
allows one to transform the integral as
\[
\Sigma_a(p^2)=-\frac{\tilde\lambda^2}{3!}\,\frac{\Gamma(2-d/2)}{(4\pi)^{\frac d2}}\,\frac{\Gamma(3-d/2)}{\Gamma(2-d/2)\,\Gamma(1)}\int_0^1du\int_0^1dw\,J(u,w)\,,
\]
\[
J(u,w)=\int\frac{d^dk_1}{(2\pi)^d}\,\frac{w^{1-\frac d2}}{\bigl[\tilde m^2+(1-w)\,k_1^2+w\,u(1-u)\,(k_1+p)^2\bigr]^{3-\frac d2}}\,.
\]
The integral over $k_1$ can be evaluated after the shift of this momentum
variable
\[
q=k_1+\frac{w\,u(1-u)}A\,p\,,
\]
where again
\[
A=(1-w)+w\,u(1-u)\,.
\]
One obtains
\[
J(u,w)=\int\frac{d^dq}{(2\pi)^d}\,\frac{w^{1-d/2}}{\Bigl(A\,q^2+w(1-w)\,u(1-u)\,p^2/A+\tilde m^2\Bigr)^{3-\frac d2}}\,.
\]
With appropriate rescaling of $q$ one can bring this integral to the form
\eqref{eq:dimregalpha}, and obtain
\[
J(u,w)=\frac{A^{3-\frac{3d}2}}{(4\pi)^{\frac d2}}\,\frac{\Gamma(3-d)}{\Gamma(3-d/2)}\,w^{1-d/2}\,\Bigl[w(1-w)\,u(1-u)\,p^2+A\,\tilde m^2\Bigr]^{d-3}\,.
\]
Combining these results we get
\[
\Sigma_a(p^2)=-\frac{\tilde\lambda^2}{3!}\,\frac{\Gamma(3-d)}{(4\pi)^d}\int_0^1dw\,du\,w^{1-\frac d2}\,A^{3-\frac{3d}2}\,\Bigl[u(1-u)\,w(1-w)\,p^2+A\,\tilde m^2\Bigr]^{d-3}\,.
\]
This quantity has second order pole in $\epsilon$, one singular factor coming
from the Gamma function
\[
\Gamma(3-d)=-\frac1\epsilon+O(1)\,,
\]
and another appearing due to the divergence of the integral over $w$ at
$w\to0$. However, subtracting $\Sigma_a(0)$ and $p^2\Sigma_a'(0)$ one obtains
\[
\Sigma_r(p^2)=-\frac{\tilde\lambda^2}{3!}\,\frac{\Gamma(3-d)}{(4\pi)^d}\,(\tilde m^2)^{d-3}\int_0^1du\,dw\,w^{1-\frac d2}\,A^{-\frac d2}\,H(u,w)\,,
\]
where
\[
H(u,w)=(1+p^2B)^{d-3}-1-(d-3)\,p^2B\,,\qquad\text{and}\qquad B=\frac{u(1-u)\,w(1-w)}{\tilde m^2\,A}\,.
\]
Note that $H(u,w=0)=0$ and therefore the $w$ integration does not produce
singularity in this expression. Moreover, one can see that $H=O(\epsilon)$,
\[
H(u,w)=\epsilon\,\Bigl[p^2B-(1+p^2B)\,\log(1+p^2B)\Bigr]+O(\epsilon^2)\,,
\]
and therefore $\Sigma_r$ has finite value at $d=4$,
\[
\Sigma_r(p^2)\Big|_{d=4}=\frac{\tilde\lambda^2}{3!}\,\frac{\tilde m^2}{(4\pi)^4}\int_0^1\frac{du\,dw}{w\,A^2}\,\Bigl[p^2B-(1+p^2B)\,\log(1+p^2B)\Bigr]
\]
At the same time the subtracted terms $\Sigma_a(0)$ and $\Sigma_a'(0)$ are
singular at $d=4$. The contributions of two diagrams (a) and (b) to $\Sigma$ then
have the form
\[
\Sigma_a(p^2)+\Sigma_b(p^2)=\Sigma_r(p^2)+\tilde\lambda^2\,\bigl(M(\epsilon)+p^2N(\epsilon)\bigr)\,,
\]
where the constants $M(\epsilon)$ and $N(\epsilon)$ have poles in $\epsilon$,
\begin{equation}\label{eq:MN}
M(\epsilon)=\frac{M_{-2}}{\epsilon^2}+\frac{M_{-1}}\epsilon+O(1)\,;\qquad
N(\epsilon)=\frac{N_{-1}}\epsilon+O(1)\,.
\end{equation}
The diagram (c) brings the contribution
\[
\Sigma_c(p^2)=\tilde\lambda^2\,\bigl(p^2z^{(2)}+b^{(2)}\bigr)\,,
\]
and our normalization condition
\[
\Sigma(0)=\Sigma'(0)=0
\]
requires the choice
\[
b^{(2)}=-M(\epsilon)\,;\qquad z^{(2)}=-N(\epsilon)
\]
of the counterterm coefficients at this order. For the proper vertex we then
have simply
\[
\Gamma^{(2)}(p^2)=p^2+\tilde m^2+\Sigma_r(p^2)+O(\tilde\lambda^3)\,.
\]

\begin{hwproblem}{4}\label{hw:4}
Within the original perturbation theory (i.e.\ with the bare mass in the
propagator $1/(p^2+m_0^2)$, and no counterterm vertices) consider the sum of all
``cactus diagrams''
\[
-\Sigma_{\text{cactus}}=
\dia[0.35]{\draw[prop,densely dashed] (-2,0) -- (2,0); \draw[prop] (0,0.6) circle (0.6); \node[vtx] at (0,0) {};}
+\dia[0.35]{\draw[prop,densely dashed] (-2,0) -- (2,0); \draw[prop] (0,0.6) circle (0.6); \draw[prop] (0,1.8) circle (0.6); \node[vtx] at (0,0) {}; \node[vtx] at (0,1.2) {};}
+\dia[0.35]{\draw[prop,densely dashed] (-2,0) -- (2,0); \draw[prop] (0,0.6) circle (0.6); \draw[prop] (-0.75,1.6) circle (0.5); \draw[prop] (0.75,1.6) circle (0.5);
  \node[vtx] at (0,0) {}; \node[vtx] at (-0.42,1.02) {}; \node[vtx] at (0.42,1.02) {};}
+\dia[0.35]{\draw[prop,densely dashed] (-2,0) -- (2.6,0); \draw[prop] (0,0.6) circle (0.6); \draw[prop] (-0.75,1.6) circle (0.5); \draw[prop] (0.75,1.6) circle (0.5);
  \draw[prop] (1.75,1.6) circle (0.5); \node[vtx] at (0,0) {}; \node[vtx] at (-0.42,1.02) {}; \node[vtx] at (0.42,1.02) {}; \node[vtx] at (1.25,1.6) {};}
+\dots
\]
and ignore all other diagrams for the mass operator. In this ``cactus
approximation'', find the relation between $m_0^2$ and $m^2$ to all orders in
$\lambda$. Use Schwinger's regularization $\Phi(p^2/\Lambda^2)=\exp(-p^2/\Lambda^2)$.
\end{hwproblem}
